%! Dividing Polynomials; Remainder and Factor Theorems — Unit 3, Lesson 3.3 — Homework
\documentclass[10pt]{article}
\usepackage{algebra2-article}
\usepackage{algebra2-boxes}
\newcommand{\TallMath}[1]{$\displaystyle #1\rule[-1.4em]{0pt}{3.2em}$}

\begin{document}

\pageheader{Unit 3, Lesson 3.3}{Homework}
\namedateperiod

\vspace{0.10in}

\begin{notesbox}{Practice}
Standard form and placeholders first, every time. When a remainder is not $0$, write the answer as
quotient $+\;\dfrac{\text{remainder}}{\text{divisor}}$.

\begin{enumerate}[leftmargin=1.9em, itemsep=11pt]
  \item \textbf{Monomial divisors.}
    \begin{enumerate}[label=(\alph*), leftmargin=1.8em, itemsep=5pt]
      \item $\dfrac{16x^5 - 24x^3 + 8x^2}{8x^2} = $ \blank{5.0cm}
      \item $\dfrac{18x^3y^2 - 27x^2y^3}{9x^2y^2} = $ \blank{5.0cm}
      \item $\dfrac{20x^4 + 5x^3 - 10x^2}{5x^2} = $ \blank{5.0cm}
    \end{enumerate}

  \item \textbf{Long division.} Show the tableau on your own paper; record the results here.
    \begin{enumerate}[label=(\alph*), leftmargin=1.8em, itemsep=5pt]
      \item $(x^3 + 2x^2 - x - 2) \div (x - 1)$: \; quotient \blank{3.0cm}, remainder \blank{1.2cm}
      \item $(2x^3 - 3x^2 + 4) \div (x - 2)$: \; quotient \blank{3.0cm}, remainder \blank{1.2cm}

        \vspace{2pt}
        What placeholder did (b) need? \blank{2.0cm}
      \item $(x^4 - 2x^3 - 7x^2 + 8x + 12) \div (x^2 - x - 6)$: \; quotient \blank{3.0cm},
        remainder \blank{1.2cm}

        \vspace{2pt}
        Why could you \emph{not} use synthetic division on (c)? \blank{5.0cm}
    \end{enumerate}

  \item \textbf{Synthetic division.} Give the quotient and the remainder.
    \begin{enumerate}[label=(\alph*), leftmargin=1.8em, itemsep=5pt]
      \item $(x^3 - 7x + 6) \div (x - 1)$: \; \blank{5.0cm}
      \item $(3x^3 + 8x^2 - x - 10) \div (x + 2)$: \; \blank{5.0cm}
      \item $(x^4 - 16) \div (x - 2)$: \; \blank{5.0cm}

        \vspace{2pt}
        In Lesson 3.2 you factored $x^4-16$ completely. Does your quotient here agree with that
        answer? Explain in one line. \blank{5.0cm}
    \end{enumerate}

  \item \textbf{Remainder Theorem --- no dividing.} Let $f(x) = x^3 - 2x^2 + 3x - 4$.
    \begin{enumerate}[label=(\alph*), leftmargin=1.8em, itemsep=5pt]
      \item Remainder on division by $(x-1)$: \blank{1.6cm} \qquad by $(x+1)$: \blank{1.6cm}
            \qquad by $(x-2)$: \blank{1.6cm}
      \item Are any of those three binomials a factor of $f$? \blank{1.4cm} \; How can you tell?
            \blank{4.6cm}
    \end{enumerate}

  \item \textbf{Factor Theorem --- test, divide, finish.} Let $p(x) = x^3 + 2x^2 - 11x - 12$.
    \begin{enumerate}[label=(\alph*), leftmargin=1.8em, itemsep=5pt]
      \item Show that $x = -1$ is a zero: $p(-1) = $ \blank{1.6cm}
      \item So \blank{1.8cm} is a factor. Divide it out: depressed polynomial \blank{3.0cm}
      \item Factor completely: $p(x) = $ \blank{5.0cm}
      \item List the zeros: \blank{4.0cm}
    \end{enumerate}

  \item \textbf{Back to a graph you already read.} On Lesson 3.0's exit ticket you read the graph of
        $h(x) = x^3 - 3x - 2$ --- it touches the $x$-axis at $x = -1$ and crosses at $x = 2$.
    \begin{enumerate}[label=(\alph*), leftmargin=1.8em, itemsep=5pt]
      \item Confirm with the Factor Theorem: $h(-1) = $ \blank{1.4cm}
      \item Divide $h(x)$ by $(x+1)$: quotient \blank{3.0cm}
      \item Factor the quotient and write $h$ completely: \blank{5.0cm}
      \item The zero $x=-1$ has multiplicity \blank{1.0cm}, which is why the graph
            \blank{2.2cm} there instead of crossing.
    \end{enumerate}

  \item \textbf{Explain.} A classmate divides $f(x)$ by $(x - 5)$ and gets a remainder of $0$. In two
        or three sentences, tell them three things they now know about $f$ --- and say which
        statement in the division statement $P = D\cdot Q + R$ makes it true. \writelines{3}
\end{enumerate}
\end{notesbox}

\begin{extensionbox}
\textbf{Run the theorems backward.}
\begin{enumerate}[label=(\alph*), leftmargin=1.8em, itemsep=5pt]
  \item Find the value of $k$ that makes $(x - 2)$ a factor of $f(x) = x^3 + kx^2 - 4x + 4$, then
        factor the resulting polynomial completely.
  \item A rectangular prism has volume $V(x) = x^3 + 6x^2 + 11x + 6$ cubic units and height
        $(x + 1)$ units. Find an expression for the area of its base, then give the two base
        dimensions as binomials.
  \item Check your answer to (b) at $x = 2$: compute $V(2)$ directly, and multiply your three
        dimensions at $x=2$. Do they agree?
  \item Synthetic division evaluates as well as it divides. Use it to find $f(3)$ for
        $f(x) = 2x^4 - 5x^3 + x - 7$, then verify by substituting $3$ directly.
\end{enumerate}
\writelines{4}
\end{extensionbox}

\begin{spiralbox}
\textbf{Coming next --- Lesson 3.4: The Rational Root Theorem.} Today you could crack any polynomial
\emph{once somebody told you a zero to test}. Tomorrow you generate that list yourself: the Rational
Root Theorem turns the constant term and the leading coefficient into a short list of candidate
zeros, synthetic division tests them, and the depressed polynomial finishes the job. Notice already
that a zero can be a fraction --- $2x^3-3x^2-11x+6$ has the zero $\tfrac12$ --- so the list will need
more than whole numbers.
\end{spiralbox}

\end{document}
